% !TeX root = ../main.tex
%%% ---------------------------------------------------------------------------
%%% 实验环境与配置
%%%
%%% 这一节对应传统实验报告的「实验仪器与材料」。
%%% 唯一的判据是：**别人照着这一节能不能复现你的结果**。
%%% 物理实验要写仪器型号、量程、精度等级；计算实验要写硬件、软件版本、
%%% 随机种子——版本号漏一个，别人就复现不出来。
%%% ---------------------------------------------------------------------------
\section{实验环境与配置}
\label{sec:setup}

\subsection{硬件与软件}
\label{sec:setup:env}

\begin{table}[htbp]
  \centering
  \caption{实验环境}
  \label{tab:env}
  \begin{tabular}{ll}
    \toprule
    {项目} & {配置} \\
    \midrule
    GPU        & NVIDIA A100-SXM4-40GB $\times$ 1 \\
    CPU        & AMD EPYC 7742，64 核 \\
    操作系统   & Ubuntu 22.04.4 LTS，内核 5.15.0 \\
    CUDA       & 12.1，cuDNN 8.9.2 \\
    Python     & 3.11.9 \\
    PyTorch    & 2.3.1 \\
    随机种子   & 0, 1, 2, 3, 4（共 5 次独立重复） \\
    \bottomrule
  \end{tabular}
\end{table}

\subsection{数据集与网络}
\label{sec:setup:data}

\dataset{} 含 \num{50000} 张训练图与 \num{10000} 张测试图，共 10 类，
分辨率 $32 \times 32$。从训练集中划出 \num{5000} 张作验证集用于选点，
测试集只在最终评估时使用一次。

网络为 ResNet-18（针对 $32\times32$ 输入把首层卷积改为 $3\times3$、
去掉最大池化）。数据增强为随机裁剪（padding \qty{4}{\pixel}）与随机水平翻转，
归一化用 \dataset{} 的通道均值与标准差。

\subsection{超参数}
\label{sec:setup:hparam}

三种优化器的公共设置：批大小 128，训练 \qty{100}{\epoch}，余弦退火，
预热 \num{5} 个 epoch。各自的特有参数见\cref{tab:hparam}。

\begin{table}[htbp]
  \centering
  \caption{各优化器的超参数}
  \label{tab:hparam}
  \begin{threeparttable}
    \begin{tabular}{l S[table-format=1.3] S[table-format=1.1] l l}
      \toprule
      {优化器} & {学习率 $\lr$} & {动量 $\mu$} & {$(\beta_1,\beta_2)$} &
      {权重衰减 $\lambda$} \\
      \midrule
      SGD-M  & 0.100 & 0.9 & {--}          & \num{5e-4} \\
      Adam   & 0.001 & {--} & $(0.9, 0.999)$ & \num{5e-4} \\
      AdamW  & 0.001 & {--} & $(0.9, 0.999)$ & \num{5e-2}\tnote{a} \\
      \bottomrule
    \end{tabular}
    \begin{tablenotes}[flushleft]
      \footnotesize
      \item[a] AdamW 的权重衰减不经过梯度，与 Adam 的 $\lambda$ 不可直接
        比较，故取值不同。这一点见\cref{eq:adamw}。
    \end{tablenotes}
  \end{threeparttable}
\end{table}
